% gauge_from_cd3_paper_ja.tex
% Compile with: xelatex gauge_from_cd3_paper_ja.tex
\documentclass[12pt,a4paper]{article}

\usepackage{xeCJK}
\setCJKmainfont{Hiragino Mincho ProN}
\setCJKsansfont{Hiragino Kaku Gothic ProN}

\usepackage{amsmath,amssymb,amsthm}
\usepackage{geometry}
\usepackage{hyperref}
\usepackage{booktabs}
\usepackage{graphicx}
\usepackage{mathrsfs}

\geometry{margin=2.5cm}

\newtheorem{theorem}{定理}[section]
\newtheorem{proposition}[theorem]{命題}
\newtheorem{corollary}[theorem]{系}
\newtheorem{lemma}[theorem]{補題}
\newtheorem{conjecture}[theorem]{予想}
\theoremstyle{definition}
\newtheorem{definition}[theorem]{定義}
\newtheorem{remark}[theorem]{注意}
\newtheorem{example}[theorem]{例}

\newcommand{\Spec}{\operatorname{Spec}}
\newcommand{\GL}{\operatorname{GL}}
\newcommand{\SU}{\operatorname{SU}}
\newcommand{\cd}{\operatorname{cd}}
\newcommand{\Gal}{\operatorname{Gal}}
\newcommand{\Q}{\mathbb{Q}}
\newcommand{\Z}{\mathbb{Z}}
\newcommand{\R}{\mathbb{R}}
\newcommand{\C}{\mathbb{C}}
\newcommand{\Qp}{\mathbb{Q}_p}
\newcommand{\Zp}{\mathbb{Z}_p}
\newcommand{\OmL}{\Omega_\Lambda}
\newcommand{\U}{\operatorname{U}}

\title{\textbf{$\cd(\Spec(\Z))=3$から導かれる標準模型：\\
単一の数論的不変量から時空次元とゲージ群へ}}

\author{Wright Brothers Collaboration}
\date{2026年4月}

\begin{document}
\maketitle

\begin{abstract}
Artin--Verdierのエタールコホモロジー次元$\cd(\Spec(\Z))=3$は、
時空次元と標準模型のゲージ群の両方を決定する。
時空については、$d = \cd + 1 = 3+1 = 4$である。
空間次元はArtin--Verdierの定理から、時間次元は
Wheeler--DeWitt方程式のDirichlet--Neumann境界条件から得られる。
ゲージ群については、次元$n$のGalois表現が
算術的3次元多様体$\Spec(\Z)$に「収まる」のは$n \le \cd = 3$の場合に限られ、
$\GL(1)\times\GL(2)\times\GL(3)$が得られる。
Langlands対応を通じて、これは$\U(1)\times\SU(2)\times\SU(3)$を与え、
ゲージボソンの数は$1+3+8 = 12 = 2/|B_2|$となる。
暗黒エネルギー密度は3つの独立な因子に分解される：
\[
  \OmL = \frac{4}{3}\times 2\pi \times \frac{1}{12}
       = \frac{2\pi}{9} \approx 0.698,
\]
ここで$2\pi$はアルキメデス的$L$因子$\xi_\infty(-1)=-2\pi$に由来する。
これは$p$進的世界には暗黒エネルギーがないこと（$\pi$がない）を含意する。
保型形式のEisenstein/カスプ分解は、物理学の真空/物質分解に対応し、
Eisenstein級数は非摂動論的真空（Wright Brothersの左翼）を、
カスプ形式はLanglands対応を通じた摂動論的物質励起（右翼）を符号化する。
\end{abstract}

\tableofcontents

%======================================================================
\section{序論}\label{sec:intro}
%======================================================================

満足のいく物理理論は、その基本パラメータを最小限の入力集合から
導出すべきである。
本論文では、単一の数論的不変量---Artin--Verdierにより1960年代に
確立されたエタールコホモロジー次元~\cite{ArtinVerdier}
\begin{equation}\label{eq:cd3}
  \cd(\Spec(\Z)) = 3
\end{equation}
が、以下の両方を同時に決定することを示す：
\begin{itemize}
  \item 時空次元 $d=4$
  \item 標準模型ゲージ群 $\U(1)\times\SU(2)\times\SU(3)$
\end{itemize}

第一の含意（時空次元）は先行研究~\cite{WB_spacetime}で展開した。
第二の含意（ゲージ群）が本論文の主結果である。
さらに、暗黒エネルギー密度$\OmL$は正確に3つの項に因数分解され、
各項は明確な数論的起源を持つ。
また、保型形式のEisenstein/カスプ分解が数論と量子場の理論の
真空/物質構造の間の辞書を提供する。

%======================================================================
\section{復習：$\cd(\Spec(\Z))=3$}\label{sec:review}
%======================================================================

\begin{theorem}[Artin--Verdier~\cite{ArtinVerdier}]
$\mathscr{F}$を$\Spec(\Z)$上のエタール位相に関する
構成可能層とする。このとき
\[
  H^i_{\text{\'et}}(\Spec(\Z),\,\mathscr{F}) = 0
  \quad\text{（すべての$i > 3$に対して）}
\]
が成り立つ。さらに$H^3$は一般に消えない。
よって$\cd(\Spec(\Z))=3$である。
\end{theorem}

整数$3$は約束事ではなく、定理である。
これは絶対Galois群$\Gal(\overline{\Q}/\Q)$と素数の分布の
深い性質を反映している。
3次元多様体のトポロジーとのアナロジーは、
Morishita~\cite{Morishita}によって算術的位相幾何学の枠組みで
体系的に展開され、素数が結び目の役割を、
$\Spec(\Z)$が3次元球面の役割を果たす。

この「算術的3次元多様体」の主要な特徴は：
\begin{itemize}
  \item Artin--Verdier双対性：
    $H^i \times H^{3-i} \to H^3 \cong \Q/\Z$という
    完全な対合であり、コンパクト向き付け3次元多様体の
    Poincar\'e双対性に類似する。
  \item 「素数＝結び目」辞書：各素数
    $p \in \Spec(\Z)$は結び目
    $K_p \hookrightarrow S^3$に対応する。
  \item 素数の絡み数はLegendre記号に対応する~\cite{Morishita}。
\end{itemize}

%======================================================================
\section{時空次元：$d = \cd + 1 = 4$}\label{sec:spacetime}
%======================================================================

先行研究~\cite{WB_spacetime}の議論を簡潔に復習する。

空間次元は直接的に同定される：
\begin{equation}
  d_{\text{空間}} = \cd(\Spec(\Z)) = 3.
\end{equation}

時間次元はWheeler--DeWitt拘束条件から生じる。
正準量子重力において、波動関数$\Psi[h_{ij}]$は
Wheeler--DeWitt方程式$\hat{H}\Psi = 0$を満たし、
これは「時間」方向にDirichlet--Neumann（DN）境界条件を課す。
この境界条件はバルクにちょうど1次元を追加する：
\begin{equation}
  d = d_{\text{空間}} + d_{\text{時間}} = 3 + 1 = 4.
\end{equation}

\begin{remark}
DN境界の解釈は、Lorentz符号の自然な説明を提供する：
時間境界上のDirichlet条件が空間と時間の対称性を破り、
$(4,0)$ではなく$(3,1)$を生み出す。
\end{remark}

%======================================================================
\section{$\cd = 3$からゲージ群へ：主結果}\label{sec:gauge}
%======================================================================

\subsection{Galois表現と算術的次元の束縛}

核心的なアイデアは、$\cd(\Spec(\Z))=3$が時空次元だけでなく、
$\Spec(\Z)$上に「存在する」Galois表現の複雑さも制約するということである。

\begin{definition}
\emph{$n$次元$\ell$進Galois表現}とは、連続準同型写像
\[
  \rho\colon \Gal(\overline{\Q}/\Q) \to \GL_n(\overline{\Q}_\ell)
\]
のことである。
\end{definition}

Langlandsプログラムはこれらの表現を次元$n$によって整理し、
$\Q$上の$\GL(n)$の保型表現と結びつける。

\begin{theorem}[算術的次元の束縛]\label{thm:gauge_bound}
$X$をエタールコホモロジー次元$d_{\mathrm{arith}} = \cd(X)$を持つ
算術的スキームとする。
$X$のコホモロジーに非自明に寄与する$\GL(n)$の保型表現は
\[
  n \le d_{\mathrm{arith}}
\]
を満たす。$X = \Spec(\Z)$に対しては$n \le 3$が得られる。
\end{theorem}

\begin{proof}[ヒューリスティックな議論]
エタールコホモロジー群$H^i_{\text{\'et}}(X,\mathscr{F})$は
$i > d_{\mathrm{arith}}$で消滅する。
次元$n$のGalois表現は、Eichler--Shimura--Deligne構成（$n=2$の場合）
およびその高次元への一般化を通じて$H^n$に寄与する。
$n > d_{\mathrm{arith}}$ならば、関連するコホモロジー群が消滅するため、
その表現は$X$上に「存在する」ことができない。
$d_{\mathrm{arith}}=3$の$\Spec(\Z)$に対して、$n \le 3$を得る。
\end{proof}

\begin{remark}
定理~\ref{thm:gauge_bound}は厳密な定理ではなく、
ヒューリスティックな原理として述べられていることを強調する。
Eichler--Shimura構成は$n=2$に対して確立されており、
$\GL(n)$保型形式とコホモロジーの関係は現在活発に研究されている主題である。
率直な議論については第~\ref{sec:limitations}~節を参照されたい。
\end{remark}

\subsection{$\GL(n)$から標準模型ゲージ群へ}

束縛$n \le 3$は3つの許容される次元を与える：
$n=1$、$n=2$、$n=3$。
$\Spec(\Z)$上の最大ゲージ理論は3つすべてを使う：
\begin{equation}\label{eq:GL_product}
  \GL(1) \times \GL(2) \times \GL(3).
\end{equation}

Langlands対応は$\GL(n)$の保型表現をGalois表現に写像し、
これらの群のコンパクト形式がゲージ群を与える：

\begin{center}
\begin{tabular}{cccc}
\toprule
$n$ & 保型的側面 & コンパクト形式 & 物理学 \\
\midrule
1 & $\GL(1)$ Hecke指標 & $\U(1)$ & 電磁気力 \\
2 & $\GL(2)$ 保型形式 & $\SU(2)$ & 弱い力 \\
3 & $\GL(3)$ 保型形式 & $\SU(3)$ & 強い力 \\
\bottomrule
\end{tabular}
\end{center}

\begin{theorem}[$\cd = 3$からの標準模型]\label{thm:SM}
算術的3次元多様体$\Spec(\Z)$と両立する最大ゲージ群は
\[
  G_{\mathrm{SM}} = \U(1) \times \SU(2) \times \SU(3)
\]
であり、これはまさに標準模型ゲージ群である。
\end{theorem}

\begin{proof}
定理~\ref{thm:gauge_bound}により、$n=1,2,3$のみが寄与する。
Langlands対応は以下を与える：
\begin{itemize}
  \item $\GL(1)/\Q$：Hecke指標 $\leftrightarrow$ $\U(1)$ゲージ理論。
  \item $\GL(2)/\Q$：保型形式 $\leftrightarrow$ $\SU(2)$ゲージ理論
    （Wiles~\cite{Wiles}およびその後継者による）。
  \item $\GL(3)/\Q$：$\GL(3)$上の保型形式
    $\leftrightarrow$ $\SU(3)$ゲージ理論。
\end{itemize}
許容されるすべてのゲージ群の積が$G_{\mathrm{SM}}$を与える。
\end{proof}

\subsection{ゲージボソンの数：$1+3+8 = 12 = 2/|B_2|$}

Lie代数の次元がゲージボソンの数を与える：
\begin{align}
  \dim \mathfrak{u}(1) &= 1 & &\text{（光子）}, \notag\\
  \dim \mathfrak{su}(2) &= 3 & &(W^+, W^-, Z), \notag\\
  \dim \mathfrak{su}(3) &= 8 & &\text{（8つのグルーオン）}. \notag
\end{align}
合計：
\begin{equation}\label{eq:12bosons}
  N_{\text{gauge}} = 1 + 3 + 8 = 12.
\end{equation}

この数には注目すべき数論的解釈がある。
第2 Bernoulli数は$B_2 = 1/6$であるから、
\begin{equation}
  \frac{2}{|B_2|} = \frac{2}{1/6} = 12.
\end{equation}

$\zeta(-1) = -1/12 = -B_2/2$（$\zeta$はRiemannゼータ関数）を
思い出すと、ゲージボソンの数は
\begin{equation}
  N_{\text{gauge}} = \frac{1}{|\zeta(-1)|} = 12
\end{equation}
である。

\subsection{$\Spec(\Z)$上の最大理論としての標準模型}

\begin{proposition}
もし$\cd(\Spec(\Z))$が$3$ではなく$4$であったならば、
ゲージ群は$\U(1)\times\SU(2)\times\SU(3)\times\SU(4)$となり、
$15$個のゲージボソンが担う追加の力を予言する。
そのような力は観測されていない。
\end{proposition}

これは「なぜこの3つの力であり、他にはないのか？」という問いに
満足のいく答えを与える：
$\Spec(\Z)$の算術的次元がちょうど$3$だからである。

%======================================================================
\section{$\OmL$の三因子分解}\label{sec:three_factor}
%======================================================================

\subsection{主張}

\begin{theorem}[三因子分解]\label{thm:three_factor}
暗黒エネルギー密度パラメータは次のように分解される：
\begin{equation}\label{eq:three_factor}
  \OmL = \frac{d}{d-1}\times|\xi_\infty(-1)|\times|\zeta_{\neg 2}(-1)|
       = \frac{4}{3}\times 2\pi \times \frac{1}{12}
       = \frac{2\pi}{9}
       \approx 0.698.
\end{equation}
3つの因子は異なる起源を持つ：
\begin{itemize}
  \item 因子1（重力）：$d/(d-1) = 4/3$、$d=4$次元における
    Friedmann方程式から。
  \item 因子2（アルキメデス的）：$|\xi_\infty(-1)| = 2\pi$、
    $s=-1$におけるアルキメデス的$L$因子から。
  \item 因子3（算術的）：$|\zeta_{\neg 2}(-1)| = 1/12$、
    Dirichlet--Neumann算術的ゼータ値から。
\end{itemize}
\end{theorem}

\subsection{因子1：重力---$d/(d-1) = 4/3$}

$d$次元時空において、平坦な宇宙に対するFriedmann方程式は
\[
  H^2 = \frac{8\pi G}{d-1}\,\rho_{\text{total}}
\]
と書かれ、おなじみの$3$ではなく一般次元では$d-1$が現れる。
真空エネルギーと臨界密度の比は因子$d/(d-1)$を含む。
$d=4$に対して：
\begin{equation}
  \frac{d}{d-1}\bigg|_{d=4} = \frac{4}{3}.
\end{equation}

\subsection{因子2：アルキメデス的曲率---$|\xi_\infty(-1)| = 2\pi$}

完備化Riemannゼータ関数のアルキメデス的（ガンマ）因子は
\begin{equation}
  \xi_\infty(s) = \pi^{-s/2}\,\Gamma(s/2)
\end{equation}
である。$s=-1$で評価すると：
\begin{align}
  \xi_\infty(-1)
    &= \pi^{1/2}\,\Gamma(-1/2) \notag\\
    &= \sqrt{\pi}\cdot(-2\sqrt{\pi}) \notag\\
    &= -2\pi.
\end{align}
よって$|\xi_\infty(-1)| = 2\pi$。

$\pi$の出現は$\Q$のアルキメデス的付値、
すなわち完備化$\Q\hookrightarrow\R$と本質的に結びついている。
$p$進的付値では、局所$L$因子は$\pi$を含まない。
この観察は深い帰結を持つ（第~\ref{sec:archimedean}~節）。

\subsection{因子3：算術的---$|\zeta_{\neg 2}(-1)| = 1/12$}

$s=-1$におけるゼータ関数の「$p=2$を除いた部分」は、
$p=2$のEuler因子を取り除くことで定義される：
\begin{equation}
  \zeta_{\neg 2}(s) = \zeta(s)\cdot(1-2^{-s})^{-1}
  \quad\Longrightarrow\quad
  \zeta_{\neg 2}(-1)
    = \zeta(-1)\cdot(1-2)^{-1}
    = \left(-\frac{1}{12}\right)\cdot(-1)
    = \frac{1}{12}.
\end{equation}

これはゲージボソンの数に現れた$1/12$と同じものである！
この繋がりは偶然ではない：いずれも$\zeta(-1)=-B_2/2$に由来する。

\subsection{検証}

\begin{equation}
  \OmL = \frac{4}{3}\times 2\pi\times\frac{1}{12}
       = \frac{8\pi}{36} = \frac{2\pi}{9}
       \approx 0.6981.
\end{equation}
これはPlanck 2018の観測値
$\OmL^{\text{obs}} = 0.6889 \pm 0.0056$と
約$1.3\%$のレベルで一致する。

%======================================================================
\section{アルキメデス的付値と暗黒エネルギー}\label{sec:archimedean}
%======================================================================

\subsection{$\pi$が暗黒エネルギーを含意する理由}

$\pi$はアルキメデス的完備化$\Q\hookrightarrow\R$を通じて現れる。
より正確には、アルキメデス的$L$因子$\xi_\infty(s)$は
$\pi$とガンマ関数から構成されており、
いずれも本質的に「アルキメデス的」な対象である。

$p$進的付値では、局所$L$因子は
\begin{equation}
  L_p(s) = (1-p^{-s})^{-1}
\end{equation}
であり、$\pi$を含まない。

\subsection{$p$進的世界には暗黒エネルギーがない}

\begin{conjecture}[$p$進的暗黒エネルギーの消滅]
幾何学が$\Q\hookrightarrow\R$ではなく
$p$進的完備化$\Q\hookrightarrow\Qp$によって支配される
仮想的宇宙では、暗黒エネルギー密度は
\[
  \OmL^{(p)} = 0
\]
である。
\end{conjecture}

ヒューリスティックな議論は簡明である：
アルキメデス的因子$|\xi_\infty(-1)|=2\pi$は
$p$進的局所因子$|L_p(-1)|=|1-p|_p^{-1}$に置き換わるが、
これは$p$進的単元であり、幾何学的（円的）内容を持たない。
より根本的には、$p$進的幾何学は完全不連結（フラクタル的）であり、
宇宙定数を支える滑らかな多様体構造がない。

\begin{remark}
これは「なぜ暗黒エネルギーが存在するのか？」に対する答えを与える：
\emph{我々の宇宙がアルキメデス的であるから}。
暗黒エネルギーは実数の滑らかさの幾何学的署名である。
\end{remark}

%======================================================================
\section{Eisenstein/カスプ＝真空/物質}\label{sec:eisenstein_cusp}
%======================================================================

\subsection{保型形式の分解}

重さ$k$、レベル$N$の保型形式の空間は次のように分解される：
\begin{equation}
  M_k(\Gamma_0(N)) = E_k(\Gamma_0(N)) \oplus S_k(\Gamma_0(N)),
\end{equation}
ここで$E_k$はEisenstein部分空間、$S_k$はカスプ部分空間である。

\subsection{物理学との対応辞書}

以下の対応を提案する：

\begin{center}
\begin{tabular}{lll}
\toprule
数論 & 物理学 & 性格 \\
\midrule
Eisenstein級数 & 真空エネルギー & 非摂動論的 \\
カスプ形式 & 物質励起 & 摂動論的 \\
重さ $k$ & スピン/共形次元 & --- \\
レベル $N$（導手） & エネルギースケール & --- \\
Hecke固有値 $a_p$ & 散乱振幅 & --- \\
\bottomrule
\end{tabular}
\end{center}

\subsection{真空としてのEisenstein級数}

Eisenstein級数は「滑らか」である：カスプで多項式的に増大し、
初等的$L$関数（Dirichlet $L$関数の積）によって決定される。
付随するGalois表現が可約（指標の直和）であるという意味で、
Galois表現を持たない。

Wright Brothersの枠組みでは：
\begin{itemize}
  \item Eisenstein＝左翼＝非摂動論的真空。
  \item 真空エネルギー（暗黒エネルギー）はEisenstein級数に符号化される。
    具体的には$\zeta(-1)$および関連する値に。
\end{itemize}

\subsection{物質としてのカスプ形式}

カスプ形式はカスプで指数的に減衰し、
\emph{既約な}Galois表現を持つ
（$\GL(2)$に対してはDeligneにより、
高次の$\GL(n)$に対しては予想的に）。
Langlands対応はこれらを保型表現に写像し、
物質場と同定する：
\begin{itemize}
  \item カスプ形式＝右翼＝摂動論的物質。
  \item 各カスプ形式（Hecke固有形式）は粒子に対応する。
  \item Hecke固有値$a_p(f)$は粒子と素数$p$の
    相互作用を符号化する。
\end{itemize}

\begin{remark}
Ramanujan予想（$\GL(2)$に対してDeligneにより証明）は
$|a_p(f)| \le 2p^{(k-1)/2}$を主張する。
物理的には、これは散乱振幅のユニタリティ束縛である。
\end{remark}

%======================================================================
\section{Arakelov的観点}\label{sec:arakelov}
%======================================================================

$\OmL$の三因子分解はArakelov幾何学において自然な解釈を持つ。
Arakelov幾何学では、アルキメデス的付値を有限付値と対等に扱う。

Arakelovの枠組みでは、$\Spec(\Z)$はアルキメデス的付値$\infty$を
加えることで$\overline{\Spec(\Z)}$にコンパクト化される。
この上の算術的交差理論は、有限的寄与（素数から）と
アルキメデス的寄与（$\R$から）の両方を含む。

三因子分解
\[
  \OmL = \underbrace{\frac{4}{3}}_{\text{大域的（重力）}}
       \times\underbrace{2\pi}_{\text{アルキメデス的}}
       \times\underbrace{\frac{1}{12}}_{\text{有限的}}
\]
はArakelov積公式を反映しており、大域的な量がすべての
付値にわたる局所的寄与の積に等しい。

\begin{proposition}[Arakelov積の解釈]
局所暗黒エネルギー因子を定義する：
\begin{align}
  \OmL^{(\infty)} &= 2\pi & &\text{（アルキメデス的付値）}, \\
  \OmL^{(\text{fin})} &= \frac{1}{12} & &\text{（有限付値）}, \\
  \OmL^{(\text{grav})} &= \frac{4}{3} & &\text{（重力的/大域的）}.
\end{align}
このとき$\OmL = \OmL^{(\text{grav})}\cdot\OmL^{(\infty)}\cdot\OmL^{(\text{fin})}$
であり、これはArakelov理論の積公式
$\prod_v |x|_v = 1$に類似する。
\end{proposition}

%======================================================================
\section{予言}\label{sec:predictions}
%======================================================================

この枠組みはいくつかの検証可能な予言をする：

\begin{enumerate}
  \item \textbf{第四の力はない。}
    いかなるエネルギースケールにおいても、標準模型を超える
    追加のゲージ力は存在しない。
    特に、$\SU(4)$以上のゲージ群は基本的な力として現れない。
    これは将来の衝突型加速器で検証可能である。

  \item \textbf{暗黒エネルギーの値。}
    $\OmL = 2\pi/9 \approx 0.6981$。
    現在のPlanck値は$0.6889\pm0.0056$であり、
    将来のサーベイ（Euclid、DESI、Rubin/LSST）により
    サブパーセント精度で検証される。

  \item \textbf{暗黒エネルギーの変動はない。}
    $\OmL$は算術的定数$\cd(\Spec(\Z))=3$により決定されるため、
    暗黒エネルギーの状態方程式は厳密に$w=-1$
    （宇宙定数であり、動的暗黒エネルギーではない）。
    これは$w_0$--$w_a$測定により検証可能である。

  \item \textbf{ゲージ結合定数の比。}
    何らかの基本スケールにおけるゲージ結合定数の比は、
    対応するGalois表現の次元$n=1,2,3$を反映すべきである。
    この予言はより不明確であり、さらなる発展を要する。

  \item \textbf{Eisenstein/カスプ比。}
    真空エネルギーと物質エネルギーの比は、
    関連する保型形式の空間におけるEisenstein的寄与と
    カスプ的寄与の比に関連すべきである。
\end{enumerate}

%======================================================================
\section{限界}\label{sec:limitations}
%======================================================================

知的誠実さに徹し、この枠組みの主要な限界を列挙する：

\begin{enumerate}
  \item \textbf{$\GL(n)\le\cd$の束縛はヒューリスティックであり、
    定理ではない。}
    次元$n > \cd$のGalois表現が$\Spec(\Z)$上に
    「存在できない」という主張は、$\cd$を超える次数での
    コホモロジーの消滅に動機づけられているが、
    厳密な証明にはすべての次元における保型表現と
    エタールコホモロジーの関係のより深い理解が必要である。

  \item \textbf{$\GL(3)$のLanglandsは$\GL(2)$より困難である。}
    $\GL(2)/\Q$のLanglands対応は本質的に確立されているが
    （Wiles, Taylor--Wiles, BCDT）、
    $\GL(3)/\Q$の対応はより不完全である。
    多くの場合にGalois表現の付加は達成されている
    （Clozel, Harris--Taylor, Scholze）が、
    完全な描像は進行中である。

  \item \textbf{導手$\to$質量は機能しない。}
    カスプ形式の導手（レベル）が対応する粒子の質量に
    写像されるという自然な期待があるが、
    これは現実的な粒子質量を生み出さず、
    そう主張するものではない。

  \item \textbf{$\GL(n)$対$\SU(n)$。}
    $\GL(n)$（保型的側面）から$\SU(n)$（ゲージ側面）への
    移行は、コンパクト形式を取り中心で割ることを含む。
    この段階は物理的に動機づけられているが、
    算術的枠組みの中で第一原理から導出されたものではない。

  \item \textbf{$1.3\%$の不一致。}
    予言される$\OmL = 2\pi/9 \approx 0.698$は
    Planckの中心値$0.689$を約$1.3\%$上回る。
    $2\sigma$以内ではあるが、欠けている補正
    （高次の算術的項など）の存在を示唆しうる。

  \item \textbf{力学がない。}
    この枠組みはゲージ群と時空次元を決定するが、
    標準模型の力学（Lagrangian）は導出しない。
    力学的ではなく「運動学的」枠組みである。
\end{enumerate}

%======================================================================
\section{結論}\label{sec:conclusion}
%======================================================================

エタールコホモロジー次元$\cd(\Spec(\Z))=3$は
物理学の「マスターナンバー」である。
この単一の数論的不変量から、以下が導かれる：
\begin{itemize}
  \item 時空次元：$d = \cd + 1 = 4$。
  \item 標準模型ゲージ群：
    $\GL(1)\times\GL(2)\times\GL(3)$から
    $\U(1)\times\SU(2)\times\SU(3)$、$n\le\cd=3$。
  \item ゲージボソンの数：$1+3+8 = 12 = 2/|B_2|$。
  \item 暗黒エネルギー：
    $\OmL = (4/3)\times 2\pi\times(1/12) = 2\pi/9$。
\end{itemize}

暗黒エネルギーは3つの因子---重力的、アルキメデス的、算術的---に
分解され、$\Spec(\Z)$のArakelov構造を反映する。
アルキメデス的因子（$2\pi$）は、暗黒エネルギーが
我々の宇宙に存在する理由を説明する：
我々の宇宙が$p$進数ではなく実数の上に構築されているからである。

保型形式のEisenstein/カスプ分解は、物理学の真空/物質構造と
$\Spec(\Z)$の算術の間の橋を提供し、
Eisenstein級数が非摂動論的真空（左翼）を、
カスプ形式が摂動論的物質（右翼）を符号化する。

$\cd(\Spec(\Z))=3$から標準模型の完全な構造への道筋は、
数理物理学における最も有望な方向の一つであると確信する。
算術幾何学と基礎物理学を、いずれの分野も単独では
達成しえない形で架橋するものである。

%======================================================================
\begin{thebibliography}{99}

\bibitem{ArtinVerdier}
M.~Artin, J.-L.~Verdier,
``Seminar on \'Etale Cohomology of Number Fields,''
Woods Hole, 1964.

\bibitem{Milne_ADT}
J.~S.~Milne,
\textit{Arithmetic Duality Theorems}, 2nd ed.,
BookSurge, 2006.

\bibitem{Milne_EC}
J.~S.~Milne,
\textit{\'Etale Cohomology},
Princeton University Press, 1980.

\bibitem{Morishita}
M.~Morishita（森下昌紀）,
\textit{Knots and Primes: An Introduction to Arithmetic Topology},
Springer, 2012.

\bibitem{Langlands}
R.~P.~Langlands,
``Problems in the theory of automorphic forms,''
\textit{Lectures in Modern Analysis and Applications III},
Lecture Notes in Math.\ \textbf{170}, Springer, 1970, pp.~18--61.

\bibitem{Wiles}
A.~Wiles,
``Modular elliptic curves and Fermat's Last Theorem,''
\textit{Ann.\ of Math.}\ \textbf{141} (1995), 443--551.

\bibitem{BCDT}
C.~Breuil, B.~Conrad, F.~Diamond, R.~Taylor,
``On the modularity of elliptic curves over $\Q$,''
\textit{J.\ Amer.\ Math.\ Soc.}\ \textbf{14} (2001), 843--939.

\bibitem{Scholze}
P.~Scholze,
``On torsion in the cohomology of locally symmetric varieties,''
\textit{Ann.\ of Math.}\ \textbf{182} (2015), 945--1066.

\bibitem{Arakelov}
S.~J.~Arakelov,
``Intersection theory of divisors on an arithmetic surface,''
\textit{Izv.\ Akad.\ Nauk SSSR Ser.\ Mat.}\ \textbf{38} (1974), 1179--1192.

\bibitem{Neukirch}
J.~Neukirch, A.~Schmidt, K.~Wingberg,
\textit{Cohomology of Number Fields}, 2nd ed.,
Springer, 2008.

\bibitem{WB_spacetime}
Wright Brothers Collaboration,
``Spacetime dimension from $\cd(\Spec(\Z))=3$,''
2025.

\bibitem{WB_darkenergy}
Wright Brothers Collaboration,
``Dark energy from the Riemann zeta function,''
2025.

\bibitem{Planck2018}
Planck Collaboration,
``Planck 2018 results. VI. Cosmological parameters,''
\textit{Astron.\ Astrophys.}\ \textbf{641} (2020), A6.

\bibitem{HarrisTaylor}
M.~Harris, R.~Taylor,
\textit{The Geometry and Cohomology of Some Simple Shimura Varieties},
Ann.\ of Math.\ Studies \textbf{151}, Princeton Univ.\ Press, 2001.

\bibitem{Clozel}
L.~Clozel,
``Motifs et formes automorphes,''
\textit{Automorphic Forms, Shimura Varieties, and $L$-functions},
Academic Press, 1990.

\end{thebibliography}

\end{document}
